
\documentclass[11pt,a4paper]{article}

\usepackage{fontspec}
\usepackage[french]{babel}
\usepackage{amsmath,amssymb,amsthm,mathtools}
\usepackage{geometry}
\usepackage{xcolor}
\usepackage{fancyhdr}
\usepackage{enumitem}
\usepackage{hyperref}
\usepackage{titlesec}
\usepackage{microtype}
\usepackage{booktabs}
\usepackage{array}

\geometry{margin=2.2cm}
\setmainfont{DejaVu Serif}
\setsansfont{DejaVu Sans}
\setmonofont{DejaVu Sans Mono}

\definecolor{accent}{HTML}{1F4E79}
\definecolor{lightaccent}{HTML}{EAF2F8}

\pagestyle{fancy}
\fancyhf{}
\lhead{\small Cours Deep Reinforcement Learning}
\chead{\small Classe : 4IA}
\rhead{\small Année : 2026}
\lfoot{\small Enseignant : Mourad Zéraï, PhD - ESPRIT School of Engineering}
\cfoot{}
\rfoot{\thepage}
\setlength{\headheight}{14pt}

\titleformat{\section}{\large\bfseries\color{accent}}{\thesection.}{0.5em}{}
\titleformat{\subsection}{\normalsize\bfseries\color{accent}}{\thesubsection.}{0.5em}{}
\titleformat{\subsubsection}{\normalsize\bfseries}{\thesubsubsection.}{0.5em}{}

\newtheorem*{definition}{Définition}
\newtheorem*{rappel}{Rappel}
\newtheorem*{hypothese}{Hypothèse}
\newtheorem*{propriete}{Propriété}
\newtheorem*{theoreme}{Théorème}
\newtheorem*{remarque}{Remarque}

\newcommand{\E}{\mathbb{E}}
\newcommand{\Pbb}{\mathbb{P}}
\newcommand{\R}{\mathbb{R}}
\newcommand{\argmax}{\operatorname*{argmax}}
\newcommand{\given}{\,|\,}

\begin{document}

\begin{center}
    {\LARGE\bfseries Dérivation et explication pas à pas de SARSA}\\[0.4em]
    {\large Note de cours rigoureuse, détaillée et appliquée à FrozenLake}
\end{center}

\vspace{0.6em}

\section{Objectif}

Le but de cette note est d'expliquer l'algorithme \emph{SARSA}, de dériver sa règle de mise à jour à partir de l'équation de Bellman d'espérance pour $q_{\pi}$, puis d'en montrer l'application à FrozenLake.

\section{Point de départ : l'équation de Bellman pour $q_{\pi}$}

Pour une politique $\pi$,
\[
q_{\pi}(s,a)=\E_{\pi}\!\left[R_{t+1}+\gamma q_{\pi}(S_{t+1},A_{t+1}) \middle\vert S_t=s,A_t=a\right].
\]
Cette équation dit que la valeur d'un couple état-action est l'espérance d'une récompense immédiate plus d'une valeur future actualisée.

\section{Passage à l'idée de \emph{bootstrap}}

Si l'on connaissait déjà $q_{\pi}$, l'équation ci-dessus permettrait de le vérifier. Mais comme $q_{\pi}$ est inconnu, SARSA remplace le terme exact
\[
q_{\pi}(S_{t+1},A_{t+1})
\]
par l'estimation courante
\[
Q(S_{t+1},A_{t+1}).
\]

\begin{remarque}
C'est ce qu'on appelle le \emph{bootstrap} : on met à jour une estimation à l'aide d'une autre estimation.
\end{remarque}

\section{Construction de la cible temporelle}

À partir d'une transition observée
\[
(S_t,A_t,R_{t+1},S_{t+1},A_{t+1}),
\]
on définit la cible temporelle
\[
Y_t = R_{t+1}+\gamma Q(S_{t+1},A_{t+1}).
\]

\paragraph{Pourquoi cette quantité est naturelle.}
L'équation de Bellman suggère que
\[
q_{\pi}(S_t,A_t)
\]
doit être proche de
\[
R_{t+1}+\gamma q_{\pi}(S_{t+1},A_{t+1}).
\]
Comme $q_{\pi}$ est inconnu, on remplace le second terme par l'estimation courante.

\section{Erreur TD}

On définit l'erreur de différence temporelle :
\[
\delta_t = Y_t - Q(S_t,A_t)
= R_{t+1}+\gamma Q(S_{t+1},A_{t+1})-Q(S_t,A_t).
\]

\paragraph{Interprétation.}
\begin{itemize}[leftmargin=1.6em]
    \item si $\delta_t>0$, la transition observée est meilleure que prévu ;
    \item si $\delta_t<0$, elle est moins bonne que prévu ;
    \item si $\delta_t=0$, l'estimation est localement cohérente avec la cible.
\end{itemize}

\section{Règle de mise à jour}

SARSA effectue la correction
\[
Q(S_t,A_t)\leftarrow Q(S_t,A_t)+\alpha \delta_t,
\]
soit explicitement
\[
Q(S_t,A_t)\leftarrow Q(S_t,A_t)+\alpha\left[R_{t+1}+\gamma Q(S_{t+1},A_{t+1})-Q(S_t,A_t)\right].
\]

\subsection{Justification élémentaire}

Cette mise à jour est une interpolation entre l'ancienne valeur et la cible :
\begin{align*}
Q_{\text{new}}(S_t,A_t)
&= Q_{\text{old}}(S_t,A_t)+\alpha\left[Y_t-Q_{\text{old}}(S_t,A_t)\right] \\
&= (1-\alpha)Q_{\text{old}}(S_t,A_t)+\alpha Y_t.
\end{align*}
Donc :
\begin{itemize}[leftmargin=1.6em]
    \item si $\alpha=1$, on remplace complètement par la cible ;
    \item si $0<\alpha<1$, on effectue une correction partielle.
\end{itemize}

\section{Pourquoi SARSA est dit \emph{on-policy}}

Le dernier terme utilise \emph{l'action réellement choisie au pas suivant}, à savoir $A_{t+1}$, qui provient de la politique courante. Si cette politique est $\varepsilon$-gloutonne, alors la mise à jour apprend la valeur de cette même politique $\varepsilon$-gloutonne.

C'est précisément la différence avec Q-learning, qui utiliserait
\[
\max_{a'}Q(S_{t+1},a')
\]
à la place de $Q(S_{t+1},A_{t+1})$.

\section{Contrôle avec SARSA}

Pour faire du contrôle, on couple la mise à jour de $Q$ avec une politique dérivée de $Q$, souvent $\varepsilon$-gloutonne.

\subsection{Politique $\varepsilon$-gloutonne}

Pour chaque état $s$ :
\[
\pi(a\given s)=
\begin{cases}
1-\varepsilon+\dfrac{\varepsilon}{|\mathcal A|} & \text{si } a\in\argmax_b Q(s,b),\\[0.8em]
\dfrac{\varepsilon}{|\mathcal A|} & \text{sinon.}
\end{cases}
\]

\subsection{Boucle logique}

\begin{enumerate}[leftmargin=1.6em]
    \item choisir $A_t$ selon la politique courante ;
    \item observer $(R_{t+1},S_{t+1})$ ;
    \item choisir $A_{t+1}$ selon la même politique ;
    \item mettre à jour $Q(S_t,A_t)$.
\end{enumerate}

\section{Application à FrozenLake}

\subsection{Transition élémentaire}

Dans FrozenLake, une transition fournit exactement les cinq objets qui donnent son nom à SARSA :
\[
(S_t,A_t,R_{t+1},S_{t+1},A_{t+1}).
\]

\subsection{Exemple concret}

Supposons qu'en état $s=0$, l'agent choisisse Right. Il obtient la transition :
\[
S_t=0,\quad A_t=\text{Right},\quad R_{t+1}=0,\quad S_{t+1}=1.
\]
Puis la politique choisit
\[
A_{t+1}=\text{Down}.
\]
La mise à jour devient
\[
Q(0,\text{Right}) \leftarrow Q(0,\text{Right}) + \alpha\left[0+\gamma Q(1,\text{Down}) - Q(0,\text{Right})\right].
\]

\subsection{Effet du mode glissant}

En mode glissant, l'action demandée peut produire une autre direction. SARSA apprend alors la valeur des actions \emph{dans l'environnement réellement subi}, y compris le risque de dérapage. Comme l'algorithme est on-policy, il tient compte du fait que la politique future continuera elle aussi à explorer.

\subsection{Conséquence pédagogique importante}

SARSA apprend souvent une politique plus prudente que Q-learning dans des environnements risqués, parce qu'il évalue la politique exploratoire réelle et non une politique idéale purement gloutonne.

\section{Pseudo-code pédagogique}

\begin{verbatim}
Initialiser Q(s,a) arbitrairement
Pour chaque épisode :
    Initialiser s
    Choisir a selon la politique epsilon-gloutonne dérivée de Q
    Répéter jusqu'à terminaison :
        Exécuter a, observer r, s'
        Choisir a' selon la même politique epsilon-gloutonne
        Q(s,a) = Q(s,a) + alpha * [r + gamma Q(s',a') - Q(s,a)]
        s = s'
        a = a'
\end{verbatim}

\section{Ce qu'il faut retenir}

\begin{itemize}[leftmargin=1.6em]
    \item SARSA part directement de l'équation de Bellman d'espérance pour $q_{\pi}$.
    \item La règle de mise à jour est obtenue en remplaçant la valeur exacte future par une estimation courante.
    \item L'erreur TD mesure le désaccord local entre prédiction et cible.
    \item SARSA est on-policy car il utilise l'action effectivement choisie au pas suivant.
    \item Dans FrozenLake glissant, cela peut conduire à des comportements plus prudents.
\end{itemize}

\end{document}
